% Source-derived chapter 10: The face $\mf$ as a product
% Original source: rigid-local-systems-multiplicative-eigenvalue-problem
\chapter{The face $\mf$ as a product}
\label{rigid-local-systems-multiplicative-eigenvalue-problem-chapter-10}
\source{rigid-local-systems-multiplicative-eigenvalue-problem}

\begin{remark}[Source section]
\label{rigid-local-systems-multiplicative-eigenvalue-problem-chapter-10-source}
\source{rigid-local-systems-multiplicative-eigenvalue-problem}
\source{rigid-local-systems-multiplicative-eigenvalue-problem}
This chapter reproduces the corresponding section of the retrieved
LaTeX source, including its definitions, statements, examples, and
proofs. It has not yet been translated into Lean.
\notready
\end{remark}

% The source section title was: The face $\mf$ as a product
We want to show the product structure \eqref{productstructure} of the face $\mf$. Recall the setting:
we have a tuple $(d,r,I^1,\dots,I^s)$ such that
$\langle\sigma_{I^1}, \sigma_{I^2},\dots,\sigma_{I^s}\rangle_{d}=1$. This defines a face $\mf$  of the monoid
$\Pic^+_{\Bbb{Q}}(\Par_{n,\mathcal{O},S})$ given by equality in inequality
\eqref{wallnew}. In Theorem \ref{Adagio} we have produced extremal ray generators $[D(a,j)]$ of $\Pic^+_{\Bbb{Q}}(\Par_{n,\mathcal{O},S})$ which lie on $\mf$.

\begin{defi}
\source{rigid-local-systems-multiplicative-eigenvalue-problem}
\notready\label{Br2}
\source{rigid-local-systems-multiplicative-eigenvalue-problem}
\hspace{2em}
\begin{enumerate}
\item Let $\Pic'=\Pic'(\mf)\subseteq\Pic(\Par_{n,\mathcal{O},S})$ be the subgroup formed by  $\mathcal{B}(\vec{\lambda},\ell)$ such that for any $j$, both of the following two conditions hold : (a) $\lambda^j_{a-1}=\lambda^{j}_{a}$ whenever $a\in I^j$, $a>1$ and $a-1\not \in I^j$ and (b) if $n\not\in I^j$ and $1\in I^j$ then $\lambda^j_1=\lambda^{j}_{n}+\ell$.
\item  Let $\Pic'^{+,\deg=0}(\mf)=\mf\cap \Pic'(\mf)$.
\end{enumerate}
\end{defi}

The product structure \eqref{productstructure} follows from the following:
\begin{proposition}
\source{rigid-local-systems-multiplicative-eigenvalue-problem}
\notready\label{Br4} Let $D_1,\dots, D_q$ be the list of elements $[D(a,j)]$ produced by Theorem \ref{Adagio}. There is an isomorphism
\source{rigid-local-systems-multiplicative-eigenvalue-problem}
$\bigoplus_{i=1}^{q}\Bbb{Q}D_i \oplus  \Pic' (\mf)_{\Bbb{Q}}\leto{\sim} \Pic^{\deg=0}_{\Bbb{Q}}(\Par_{n,\mathcal{O},S})$, where $\Pic^{\deg=0}_{\Bbb{Q}}(\Par_{n,\mathcal{O},S})$ is the subset of $\Pic(\Par_{n,\mathcal{O},S})$ formed by line bundles for which equality holds in
\eqref{wallnew}.
Setting $\mf^{(2)}= \Pic'^{+,\deg=0}_{\Bbb{Q}}(\mf) $, we obtain a bijection
$$\prod_{i=1}^{q}\Bbb{Q}_{\geq 0}D_i \times  \mf^{(2)} \leto{\sim} \mf.$$
\end{proposition}




\begin{lemma}
\source{rigid-local-systems-multiplicative-eigenvalue-problem}
\notready\label{oneoff}
\source{rigid-local-systems-multiplicative-eigenvalue-problem}
Each of the divisors $D(a,j)$ in Theorem \ref{Adagio}, fails exactly one of the equalities defining $\Pic'$:
Write $\mathcal{O}(D(a,j))=\mathcal{B}(\vec{\lambda},\ell)$.
\begin{enumerate}
\item If $a>1$, then $\lambda^j_{a-1}=\lambda^j_{a}+1$.
\item If $a=1$, then $\lambda^j_n =\lambda^j_1-\ell +1$.
\end{enumerate}
\end{lemma}
\begin{proof}
These statements follow from Theorem \ref{brahms} and  equation \eqref{ABL4}.
\end{proof}

\subsection{Proof of Proposition \ref{Br4}}

The proof is a direct generalization of \cite[Theorem 1.15]{BHermit}. The main point is the following: The first isomorphism follows immediately from Lemma \ref{oneoff}. For the second bijection we proceed as follows. Suppose we have  a point $\mathcal{B}(\vec{\lambda},\ell)\in\mf$ not in $\Pic'$. Assume first that it fails an inequality $\lambda^j_{a-1}=\lambda^{j}_{a}$ with $a\in I^j$, $a>1$ and $a-1\not\in I^i$. We will then show that any global section of  $\mathcal{B}(\vec{\lambda},\ell)$ vanishes at any point of the divisor $D(a,j)$, so that $\mathcal{B}(\vec{\lambda},\ell)(-D(a,j))$ is still effective. It still lies on $\mf$ because $D(a,j)$ is on $\mf$ by Theorem \ref{Adagio}.

To show the vanishing, let $(\mv,\mw,\me,\gamma)$ be a general point of $D(a,j)$. The semistability inequality \eqref{wallnew} holds as an equality for $\mathcal{B}(\vec{\lambda},\ell)$. The semistability inequality corresponding to $\mv$  for $\mathcal{B}(\vec{\lambda},\ell)$  replaces $\lambda^j_a$ in the above equality by $\lambda^j_{a-1}$ and hence fails because of Lemma \ref{oneoff}. Therefore any global section of $\mathcal{B}(\vec{\lambda},\ell)$ vanishes on $D(a,j)$ as desired.


If $\mathcal{B}(\vec{\lambda},\ell)$  fails an inequality  $\lambda^j_1=\lambda^{j}_{n}+\ell$ with $n\not\in I^j$ and $1\in I^j$ then we proceed similarly: In this case $\mv$ is only a subsheaf of $\mw$, but we can consider the subbundle given by $\mv$ (``the saturation''). The semistability inequality corresponding to the saturation is violated at a generic point of $D(1,j)$: In this case the left hand side of $\eqref{wallnew}$ decreases by $\ell-1$ and the right hand side by $\ell$ since $d$ drops by one (by Lemma \ref{oneoff}).


By repeatedly carrying out these subtractions we can write any element of  $\mf$ as a sum corresponding to the second bijection.
\begin{remark}
\source{rigid-local-systems-multiplicative-eigenvalue-problem}
\notready
As in \cite{BHermit,BKiers,Kiers}, we may distinguish two types of rays of $\mathcal{F}$. A ray
$\Bbb{Q}_{\geq 0}(\vec{\lambda},\ell)$ is a type I ray of $\mathcal{F}$ if $\mathcal{B}(\vec{\lambda},\ell)$ is not in $\Pic'(\mf)$, and type II otherwise. The extremal rays  generated by $D(a,j)$  are  type I on $\mf$. An extremal ray of  $\Pic^+_{\Bbb{Q}}(\Par_{n,\mathcal{O},S})$ may lie on two different regular facets. It is possible that it is a type I ray of one facet and a type II of the other. All rays that are type I on some regular facet are generated by F-line bundles.
\end{remark}
